% Generated draft: requires mathematical and source review.
\documentclass[11pt]{article}
\usepackage[margin=25mm]{geometry}
\usepackage{fontspec}
\setmainfont{Latin Modern Roman}
\usepackage{amsmath,amssymb,mathtools,unicode-math}
\setmathfont{Latin Modern Math}
\usepackage{enumitem,needspace,xurl,hyperref}
\setlength{\parindent}{0pt}
\setlength{\parskip}{5pt}
\setlength{\emergencystretch}{3em}
\newcommand{\N}{\mathbb{N}}
\newcommand{\Z}{\mathbb{Z}}
\newcommand{\Q}{\mathbb{Q}}
\newcommand{\R}{\mathbb{R}}
\newcommand{\C}{\mathbb{C}}
\newcommand{\F}{\mathbb{F}}
\newcommand{\A}{\mathbb{A}}
\newcommand{\PP}{\mathbb{P}}
\newcommand{\E}{\mathbb{E}}
\newcommand{\eps}{\varepsilon}
\newcommand{\dd}{\,\mathrm{d}}
\begin{document}
\section*{UnsolvedMath ID: 800373}
\textbf{Source problem code:} catalogue entry
\subsection{Definitions and mathematical statement}
Let \(d\geq 0\) be an integer and let
\[
H(t)=\sum_{i=0}^{d}h_i t^i\in\Z_{\geq 0}[t]
\]
be a polynomial of degree \(d\). Assume that \(H\) is symmetric, meaning \(h_i=h_{d-i}\) for every \(0\leq i\leq d\), monic, meaning \(h_d=1\), and real-rooted, meaning that every complex zero of \(H\) belongs to \(\R\).

A simplicial convex \(d\)-polytope is a convex polytope of dimension \(d\) whose nonempty proper faces are simplices. For such a polytope \(P\), let \(f_{j-1}(P)\) denote the number of its \((j-1)\)-dimensional faces for \(1\leq j\leq d\), and put \(f_{-1}(P)=1\). Its \(h\)-vector \((h_0(P),\ldots,h_d(P))\) is defined by
\[
\sum_{i=0}^{d}h_i(P)t^i
=
\sum_{j=0}^{d}f_{j-1}(P)t^j(1-t)^{d-j}.
\]

\textbf{Question (Brenti, Problem 11).} For every polynomial \(H\) satisfying the preceding hypotheses, does there exist a simplicial convex \(d\)-polytope \(P\) such that
\[
h_i(P)=h_i\qquad(0\leq i\leq d)?
\]
The source does not explicitly specify a lower bound on \(d\) or a zero-dimensional convention. If \(d=0\) is included, the displayed definition uses the convention that a zero-dimensional polytope has \(h\)-vector \((1)\).
\subsection{Short English statement}
Brenti asks whether every symmetric, monic polynomial with nonnegative integer coefficients and only real roots is the \(h\)-polynomial of a simplicial convex polytope of the corresponding dimension. The supplied registry labels the question open; that label is not an independently verified literature assessment.
\subsection{Sources}
\begin{itemize}
\item Supplied local UnsolvedMath registry record, permanent ID \(800373\), source problem code \texttt{catalogue entry}, title ``Coxeter Groups and Unimodality --- Problem 11.'' The record attributes the question to Brenti, with proposed year 2024, and labels its status ``open.'' Upstream revision: \texttt{d371ee917ac9647fbb06cd6fc529deec6f3dda82}.
\item Brenti, \emph{Some Open Problems on Coxeter Groups and Unimodality} (2024), Problem 11. Supplied source URL: \url{https://arxiv.org/abs/2410.09897}. The supplied provenance records access on 2026-10-16 and extraction by \texttt{arxiv\_latex}; these are supplied provenance claims, not independently verified access or source checks.
\end{itemize}
\end{document}
